\documentclass[10pt,a4paper]{amsart}
\usepackage{amsmath,amssymb,amsthm,graphicx}
% Retained minimal preamble for the Stacks Project chapter sites-cohomology.tex (Cohomology on Sites).
% Verbatim excerpt of the chapter preamble: Xy-pic setup (native preamble lines 14-18) and the native
% theorem declarations with the mathematical macros (native preamble lines 154-217). Copied mechanically;
% nothing added, removed or redefined. The wrapper class and the amsmath/amssymb/amsthm/graphicx base
% packages are supplied by the assembler, so the class/externaldocument part of the native preamble is omitted.
% For commutative diagrams we use Xy-pic
\usepackage[all]{xy}

% We use 2cell for 2-commutative diagrams.
\xyoption{2cell}
% Theorem environments.
%
\theoremstyle{plain}
\newtheorem{theorem}[subsection]{Theorem}
\newtheorem{proposition}[subsection]{Proposition}
\newtheorem{lemma}[subsection]{Lemma}

\theoremstyle{definition}
\newtheorem{definition}[subsection]{Definition}
\newtheorem{example}[subsection]{Example}
\newtheorem{exercise}[subsection]{Exercise}
\newtheorem{situation}[subsection]{Situation}

\theoremstyle{remark}
\newtheorem{remark}[subsection]{Remark}
\newtheorem{remarks}[subsection]{Remarks}

\numberwithin{equation}{subsection}

% Macros
%
\def\lim{\mathop{\mathrm{lim}}\nolimits}
\def\colim{\mathop{\mathrm{colim}}\nolimits}
\def\Spec{\mathop{\mathrm{Spec}}}
\def\Hom{\mathop{\mathrm{Hom}}\nolimits}
\def\Ext{\mathop{\mathrm{Ext}}\nolimits}
\def\SheafHom{\mathop{\mathcal{H}\!\mathit{om}}\nolimits}
\def\SheafExt{\mathop{\mathcal{E}\!\mathit{xt}}\nolimits}
\def\Sch{\mathit{Sch}}
\def\Mor{\mathop{\mathrm{Mor}}\nolimits}
\def\Ob{\mathop{\mathrm{Ob}}\nolimits}
\def\Sh{\mathop{\mathit{Sh}}\nolimits}
\def\NL{\mathop{N\!L}\nolimits}
\def\CH{\mathop{\mathrm{CH}}\nolimits}
\def\proetale{{pro\text{-}\acute{e}tale}}
\def\etale{{\acute{e}tale}}
\def\QCoh{\mathit{QCoh}}
\def\Ker{\mathop{\mathrm{Ker}}}
\def\Im{\mathop{\mathrm{Im}}}
\def\Coker{\mathop{\mathrm{Coker}}}
\def\Coim{\mathop{\mathrm{Coim}}}

% Boxtimes
%
\DeclareMathSymbol{\boxtimes}{\mathbin}{AMSa}{"02}

%
% Macros for moduli stacks/spaces
%
\def\QCohstack{\mathcal{QC}\!\mathit{oh}}
\def\Cohstack{\mathcal{C}\!\mathit{oh}}
\def\Spacesstack{\mathcal{S}\!\mathit{paces}}
\def\Quotfunctor{\mathrm{Quot}}
\def\Hilbfunctor{\mathrm{Hilb}}
\def\Curvesstack{\mathcal{C}\!\mathit{urves}}
\def\Polarizedstack{\mathcal{P}\!\mathit{olarized}}
\def\Complexesstack{\mathcal{C}\!\mathit{omplexes}}
% \Pic is the operator that assigns to X its picard group, usage \Pic(X)
% \Picardstack_{X/B} denotes the Picard stack of X over B
% \Picardfunctor_{X/B} denotes the Picard functor of X over B
\def\Pic{\mathop{\mathrm{Pic}}\nolimits}
\def\Picardstack{\mathcal{P}\!\mathit{ic}}
\def\Picardfunctor{\mathrm{Pic}}
\def\Deformationcategory{\mathcal{D}\!\mathit{ef}}
\usepackage{cleveref}
\crefname{originallemma}{original lemma}{original lemmas}
\Crefname{originallemma}{Original lemma}{Original lemmas}
\makeatletter
\newlabel{derived-lemma-injective-resolutions-exist}{{013K}{source}{}{sourcepointer-1}{}}
\newlabel{derived-lemma-injective-resolutions-exist@cref}{{[originallemma][0][]013K}{[1][0][]source}{}{}{}}
\makeatother
\crefname{originallemma}{original lemma}{original lemmas}
\Crefname{originallemma}{Original lemma}{Original lemmas}
\makeatletter
\newlabel{derived-lemma-leray-acyclicity}{{015E}{source}{}{sourcepointer-2}{}}
\newlabel{derived-lemma-leray-acyclicity@cref}{{[originallemma][0][]015E}{[1][0][]source}{}{}{}}
\makeatother
\crefname{originallemma}{original lemma}{original lemmas}
\Crefname{originallemma}{Original lemma}{Original lemmas}
\makeatletter
\newlabel{injectives-lemma-derived-products}{{07D9}{source}{}{sourcepointer-3}{}}
\newlabel{injectives-lemma-derived-products@cref}{{[originallemma][0][]07D9}{[1][0][]source}{}{}{}}
\makeatother
\crefname{originalequation}{original equation}{original equations}
\Crefname{originalequation}{Original equation}{Original equations}
\makeatletter
\newlabel{sites-modules-equation-map-lower-shriek-OU-into-module}{{0G1V}{source}{}{sourcepointer-4}{}}
\newlabel{sites-modules-equation-map-lower-shriek-OU-into-module@cref}{{[originalequation][0][]0G1V}{[1][0][]source}{}{}{}}
\makeatother
\title{If the cohomology functors on an object commute with direct sums then the extension by zero of the structure sheaf is compact}
\author{The Stacks Project Authors}
\date{stacks-project:sites-cohomology.tex:local-snapshot-20261011; GNU FDL 1.2 or later}
\begin{document}
\maketitle
\noindent\textit{Editorial note.} All delivered mathematical text is the author's own, copied from the pinned Stacks Project chapter file sites-cohomology.tex; no statement or proof step was rewritten, reordered or removed.\par\medskip
\noindent\textit{Editorial note.} The counted proof is the author's complete argument: the proof is written out in full: the reduction of the compactness statement to commutation of H\textasciicircum{}p(U, -) with direct sums, the acyclicity of direct sums of injectives, and the resolution argument for a bounded below object.\par\medskip
\noindent\textit{Editorial note.} The label lemma-when-jshriek-lower-compact names the counted statement, and the proof environment immediately following it in the source is delivered verbatim as its proof.\par\medskip
\noindent\textit{Editorial note.} The counted statement is carried by the author's own lemma environment in the chapter source.\par\medskip
\noindent\textit{Editorial note.} The statement and proof are copied from the pinned chapter file; the four cited inputs are kept as external pointers with their canonical Stacks tags.\par\medskip
\noindent\textit{Editorial note.} The preamble is the chapter's own retained minimal preamble (sources/stacks-sites-cohomology/retained-preamble.tex), a verbatim excerpt of the native preamble: the Xy-pic setup and the native theorem declarations with the mathematical macros.\par\medskip
\noindent\textit{Allowed standard prerequisites (editorial).}\begin{itemize}
\item The identification of Hom(j\_\{U!\}O\_U, -) with evaluation at U. Modules on Sites, the equation describing maps out of the extension by zero of the structure sheaf. Supplies the reduction of the compactness statement to the commutation of H\textasciicircum{}p(U, -) with direct sums.
\item Injective resolutions and acyclicity of injectives. Derived Categories, the existence of injective resolutions and the acyclicity lemma for bounded below complexes of injectives. Supplies the resolution used to compute H\textasciicircum{}p(U, M) after the reduction to bounded below M.
\item Derived products of injectives. Injectives, the lemma allowing products of injective objects to be used in the resolution argument. Supplies the control of the direct sums of injectives in the computation.
\end{itemize}\medskip
\ifdefined\hypertarget\hypertarget{sourcepointer-1}{}\fi
\noindent\textit{Source reference.} Original lemma 013K: sources/stacks-chow/tags.txt, source lines 1438--1438. Prerequisite pointer invoked at proof line(s) 14094; it names a standard fact of a sibling chapter that the counted argument uses, and every step of the counted proof built on it is retained verbatim in the delivered body.\par\medskip
\ifdefined\hypertarget\hypertarget{sourcepointer-2}{}\fi
\noindent\textit{Source reference.} Original lemma 015E: sources/stacks-chow/tags.txt, source lines 1505--1505. Prerequisite pointer invoked at proof line(s) 14098; it names a standard fact of a sibling chapter that the counted argument uses, and every step of the counted proof built on it is retained verbatim in the delivered body.\par\medskip
\ifdefined\hypertarget\hypertarget{sourcepointer-3}{}\fi
\noindent\textit{Source reference.} Original lemma 07D9: sources/stacks-chow/tags.txt, source lines 9404--9404. Prerequisite pointer invoked at proof line(s) 14095; it names a standard fact of a sibling chapter that the counted argument uses, and every step of the counted proof built on it is retained verbatim in the delivered body.\par\medskip
\ifdefined\hypertarget\hypertarget{sourcepointer-4}{}\fi
\noindent\textit{Source reference.} Original equation 0G1V: sources/stacks-chow/tags.txt, source lines 20035--20035. Prerequisite pointer invoked at proof line(s) 14081; it names a standard fact of a sibling chapter that the counted argument uses, and every step of the counted proof built on it is retained verbatim in the delivered body.\par\medskip
\begin{lemma}
\label{lemma-when-jshriek-lower-compact}
Let $(\mathcal{C}, \mathcal{O})$ be a ringed site. Let $U$ be an object of
$\mathcal{C}$. Assume the functors $\mathcal{F} \mapsto H^p(U, \mathcal{F})$
commute with direct sums. Then $\mathcal{O}$-module $j_!\mathcal{O}_U$ is a
compact object of $D^+(\mathcal{O})$ in the following sense:
if $M = \bigoplus_{i \in I} M_i$ in $D(\mathcal{O})$ is
bounded below, then $\Hom(j_{U!}\mathcal{O}_U, M) =
\bigoplus_{i \in I} \Hom(j_{U!}\mathcal{O}_U, M_i)$.
\end{lemma}

\begin{proof}
Since $\Hom(j_{U!}\mathcal{O}_U, -)$ is the same as the functor
$\mathcal{F} \mapsto \mathcal{F}(U)$ by
Modules on Sites, Equation
(\ref{sites-modules-equation-map-lower-shriek-OU-into-module}) it suffices
to prove that $H^p(U, M) = \bigoplus H^p(U, M_i)$.
Let $\mathcal{I}_i$, $i \in I$ be a collection of injective
$\mathcal{O}$-modules. By assumption we have
$$
H^p(U, \bigoplus\nolimits_{i \in I} \mathcal{I}_i) =
\bigoplus\nolimits_{i \in I} H^p(U, \mathcal{I}_i) = 0
$$
for all $p$. Since $M = \bigoplus M_i$ is bounded below, we
see that there exists an $a \in \mathbf{Z}$ such that $H^n(M_i) = 0$
for $n < a$. Thus we can choose complexes of injective $\mathcal{O}$-modues
$\mathcal{I}_i^\bullet$ representing $M_i$
with $\mathcal{I}_i^n = 0$ for $n < a$, see
Derived Categories, Lemma \ref{derived-lemma-injective-resolutions-exist}.
By Injectives, Lemma \ref{injectives-lemma-derived-products}
we see that the direct sum complex $\bigoplus \mathcal{I}_i^\bullet$
represents $M$. By Leray acyclicity
(Derived Categories, Lemma \ref{derived-lemma-leray-acyclicity})
we see that
$$
R\Gamma(U, M) = \Gamma(U, \bigoplus \mathcal{I}_i^\bullet) =
\bigoplus \Gamma(U, \bigoplus \mathcal{I}_i^\bullet) =
\bigoplus R\Gamma(U, M_i)
$$
as desired.
\end{proof}

\end{document}
